\section{
    Conclusão
} \label{sec:conclusao}

{\color{red} REVISAR TODO MUNDO JUNTO}

Este trabalho demonstrou que a segurança em aplicações web consolidou-se como um pilar estratégico, transcendendo a camada perimetral de rede \cite{CloudSecurityAlliance2024TopThreats, Alwan2021OWASPTopTenSurvey}. A defesa em profundidade surge como o princípio unificador fundamental, integrando a análise de código (SAST/IAST), testes dinâmicos (DAST) e proteção em \textit{runtime} (RASP) para mitigar falhas desde o \textit{design} até a execução \cite{DeJesus2023SecurityTestingWebApps, Alsaedi2023DeepLearningWAF, Barth2008RobustDefensesCSRF, Pahl2016MicroservicesSurvey}.

Apesar dos avanços, identificou-se uma lacuna na literatura e na prática atual quanto à segurança de arquiteturas descentralizadas, como \textit{shadow} APIs e o processamento complexo no lado do cliente (\textit{DOM-based}), que dificultam o rastreamento automatizado de vulnerabilidades \cite{Escape2023GraphQLSecurity, Tanveer2025RESTfulAPISecurity, SuredaRiera2025TwentyTwoYearsXSS}. Além disso, embora a Inteligência Artificial (IA) seja promissora na detecção de anomalias, sua falta de interpretabilidade ainda limita a auditoria crítica em sistemas de larga escala \cite{Yadav2025CuttingEdgeIDS, Alsaedi2023DeepLearningWAF}.

Como direções futuras, aponta-se a necessidade de amadurecer modelos de IA explicável aplicados à cibersegurança e a consolidação de arquiteturas \textit{Zero Trust} baseadas em identidades dinâmicas \cite{Yadav2025CuttingEdgeIDS, NIST2020ZeroTrustArchitecture, Yadav2025ZeroTrustMicroservices}. A evolução para uma cultura \textit{DevSecOps} plenamente automatizada é essencial para garantir que a agilidade no desenvolvimento não comprometa a integridade dos dados e a resiliência das infraestruturas modernas \cite{Edmundson2022SANSDevSecOps, DeJesus2023SecurityTestingWebApps}.